\documentclass[12pt]{report}

\usepackage{amssymb,amsthm}
\usepackage{amsmath,color}
%\usepackage{dsfont}
\usepackage{setspace}
\usepackage{graphicx}
\usepackage{multicol}
\usepackage[document]{ragged2e}
\usepackage{fancyhdr}
% \usepackage[]{classicthesis}
\usepackage{hyperref}
\hypersetup{
  colorlinks   = true,    % Colours links instead of ugly boxes
  urlcolor     = blue,    % Colour for external hyperlinks
  linkcolor    = blue,    % Colour of internal links
  citecolor    = red      % Colour of citations
  }
% \usepackage{cref}
  
\def\changemargin#1#2{\list{}{\rightmargin#2\leftmargin#1}\item[]}
\let\endchangemargin=\endlist 

\def\w{\mathtt w}
\def\K{\mathcal K}
\def\N{\mathbb N}
\def\R{\mathcal R}
\def\l{\ell}
\def\S{\mathcal S}
\def\Z{\mathbb Z}
\def\ora{\overrightarrow}


\newtheorem{theorem}{Theorem}[section]
\newtheorem{corollary}{Corollary}[theorem]
\newtheorem{lemma}[theorem]{Lemma}


\begin{document}

\begin{center}
\LARGE{\textbf{Indian Institute of Technology, Delhi}}\\
\vspace{0.8cm}
\large{\textbf{Design Practices in Computer Science}}\\[5pt]
\large{\textbf{COP290}}\\[5pt]

\vspace{0.5cm}

\large{\textbf{Starling Simulation }}

\begin{center}
\includegraphics[height=5cm]{iitd.eps}
\end{center}
\vspace{0.2cm}

\textbf{April 18, 2018} \\
\textbf{Department of Computer Science and Engineering} \\
\textbf{Indian Institute of Technology, Delhi}\\


\vspace{1.5cm}

\begin{multicols*}{2}

\begin{flushleft}

\textit{Authors :\\ }


\textbf{Udit Jain} \\
(2016CS10327)\\
\textbf{Shashank Goel}
(2016CS10332)\\

\end{flushleft}


\columnbreak

\begin{flushleft}

\textit{\\Supervisor :\\ }
\textbf{Prof. Subhashish Bannerjee} \\[5pt]

\end{flushleft}

\end{multicols*}

\end{center}

\newpage

\begin{center}
\Large \bf ABSTRACT
\end{center}
\vspace{0.2in}

We are going to design and implement a Starling Simulation that will help us understand the Flocking behaviour of birds. Computer Science term for this simulation is Boids.
\\
\vspace{0.3cm}
The simulation software will have the following functionalities:

\begin{enumerate}
  \item
  We will be able to interactively input or read from a file :
  \begin{itemize}
    \item
    Degree of Seperation, Cohesion, Alignment.  
    \item
    Speed of boids .
  \end{itemize}
  \item
  We can interactively add / remove \textit{ boids and obstacles } based on combination of inputs from keyboard and mouse. 
  \item
  Get the color of boids based by grouping on speed, flock, direction .
  \item
  Calcuate and display average Velocity and Accleration vector of boid groups and global flock dynamically.
  \item
  Calcuate and display average Kinetic Energy of boid groups and global flock dynamically.
\end{enumerate}


\vspace{0.5cm}
Future scope of the problem may include:

\begin{enumerate}
  \item
  Interactively control the Gravity vector and see it's effect on the flock .
  \item
  Adding and measuring the effect of Wind speed / direction and Drag on the flock .
\end{enumerate}


\vspace{0.5cm}

In this design project, we shall work as observers, developers and algorithm enthusiasts to understand the ways and finding different means to approach and tackle the objectives in a more well defined mathematical way. 

\vspace{0.2cm}
We shall work with full confidence and zeal to achieve the goal or reach to quite an end of the problem so that using our lemmas, proofs and knowledge, someday a perfect model can be implemented using a software by some other Computer Explorer. As a matter of interest, we just wish to argue that these things can be computed by our brain so we do hope to find a solution to this problem using machine learning algorithms. Since, Machine Learning algorithms are more or less based on Mathematical matrices, with the use of computer graphics, we expect to find a start with matrices that we have dealt with further in this report. 

\newpage

\tableofcontents

\newpage

% \setcounter{secnumdepth}{-1} % to unNumber introduction

\chapter{Introduction}

Boids is an artificial life program, developed by Craig Reynolds in 1986,
which simulates the flocking behaviour of birds. His paper on this topic
was published in 1987 in the proceedings of the ACM SIGGRAPH con-
ference. The name ”boid” corresponds to a shortened version of ”bird-oid
object”, which refers to a bird-like object.

\vspace{1cm}
As with most artificial life simulations, Boids is an example of emergent
behavior; that is, the complexity of Boids arises from the interaction of individual agents (the boids, in this case) adhering to a set of simple rules.

\hyperref[sec:baseline]{Baseline Rules}
% \cref{sec:baseline}
% {Basic Rules}

\vspace{0.5cm}

The following objectives are aimed to be discussed in this paper:
\begin{itemize}
  \item
  Defining the problem.
  % Support you design decisions ? 
  \item
  Discussing how a simulation is modeled, what are the diffrent padigrams.
  % Discret even simulation, Mutex and Tsigma force . screen Process Interaction, Interprocess Simulation, 
  % logical time (when equal to) real time, for loop(not giving perefrence to one starlilng) concurrency (Showing it is equal because all threads have to access each other's memory, locking problem ?) ? repeated ? ||| Bound on number of starlings ? Experimenation with hyperparameters ?
  % Changing according to some hyper-parameters ? 
  \item
  What set of rules govern bird in flocks.
  % Power, drag, float force, Can birds sustain that ? Do we limit 
  \item
  What model is to be adopted for interplay of starlings with the environment.
  \item
  What parameters can be varied to see interesting patterns.
  % (dynamic and average) power, KE . Energy expenditure, food ? , seperation radius, alignment etc
  \item
  Intresting Obeservations on varying diffrent parameters in simulation.
  % Lone bird cannot reach max speed
  \item 
  Processing Software
  % How does it work ? How does it handle threading ?
\end{itemize}

\chapter{Defining the problem}
\section{Introduction}
\subsection{Why do we care about bird flocking patterns ?}
\section{How do we study them ?}

\chapter{Modelling a Simulation}
\section{Introduction}
Simulation is the imitation of the operation of a real-world process or system. The act of simulating something first requires that a model be developed; this model represents the key characteristics, behaviors and functions of the selected physical or abstract system or process. The model represents the system itself, whereas the simulation represents the operation of the system over time.
\\
Simulation is used in many contexts, such as simulation of technology for performance optimization, safety engineering, testing, training, education, and video games. Often, computer experiments are used to study simulation models. Simulation is also used with scientific modelling of natural systems or human systems to gain insight into their functioning, as in economics. Simulation can be used to show the eventual real effects of alternative conditions and courses of action. Simulation is also used when the real system cannot be engaged, because it may not be accessible, or it may be dangerous or unacceptable to engage, or it is being designed but not yet built, or it may simply not exist.
\\
Key issues in simulation include acquisition of valid source information about the relevant selection of key characteristics and behaviours, the use of simplifying approximations and assumptions within the simulation, and fidelity and validity of the simulation outcomes. Procedures and protocols for model verification and validation are an ongoing field of academic study, refinement, research and development in simulations technology or practice, particularly in the field of computer simulation.

\subsection{What's a simulation ?}

A computer simulation (or "sim") is an attempt to model a real-life or hypothetical situation on a computer so that it can be studied to see how the system works. By changing variables in the simulation, predictions may be made about the behaviour of the system. It is a tool to virtually investigate the behaviour of the system under study.
\\
Computer simulation has become a useful part of modeling many natural systems in physics, chemistry and biology, and human systems in economics and social science (e.g., computational sociology) as well as in engineering to gain insight into the operation of those systems. A good example of the usefulness of using computers to simulate can be found in the field of network traffic simulation. In such simulations, the model behaviour will change each simulation according to the set of initial parameters assumed for the environment.
\\
Traditionally, the formal modeling of systems has been via a mathematical model, which attempts to find analytical solutions enabling the prediction of the behaviour of the system from a set of parameters and initial conditions. Computer simulation is often used as an adjunct to, or substitution for, modeling systems for which simple closed form analytic solutions are not possible. There are many different types of computer simulation, the common feature they all share is the attempt to generate a sample of representative scenarios for a model in which a complete enumeration of all possible states would be prohibitive or impossible.
\\
Several software packages exist for running computer-based simulation modeling (e.g. Monte Carlo simulation, stochastic modeling, multimethod modeling) that makes all the modeling almost effortless.
\\
Modern usage of the term "computer simulation" may encompass virtually any computer-based representation.

\section{Types of Simulation}
\subsection{Key Aspects}

\chapter{Flocking Rules}
\section{Baseline}
\label{sec:baseline}

The rules applied in the simplest Boids world are as follows:

\begin{itemize}
  \item
  \textbf{\textit{Seperation}} : Steer to avoid crowding of local flockmates
    \begin{center}
      \includegraphics[height=5cm]{Rule_separation.png}
    \end{center}
  \item
  \textbf{\textit{Alignment}} : Steer towards the average heading of local flockmates.
  \begin{center}
    \includegraphics[height=5cm]{Rule_alignment.png}
  \end{center}
  \item
  \textbf{\textit{Cohesion}} : Steer to move toward the average position (Center of Mass) of local flockmates.
    \begin{center}
    \includegraphics[height=5cm]{Rule_cohesion.png}
  \end{center}
\end{itemize}

\section{Additional Behaviour}

More complex rules can be added, such as obstacle avoidance and goal seeking .\\

\chapter{Environment Parameters}
\section{Parameters}
Press H for the full helpful Legend.
\\
Wind speed,
\section{Observations}

\chapter{Processing Software}
\section{How does it work?}
How does it work ? How does it handle threading ?
Explicit / Implicit / Automatic multi threading ? Arduion loops ?

\chapter{Conclusion}


% \begin{thebibliography}{9}


% \bibitem{ALM}
% P. Allen, B. Landman and H. Meeks, New Bounds on van der Waerden type numbers for Generalized $3$-term Arithmetic Progressions, {\it arXiv: 1201.3842v2}


% \bibitem{BL}
% S. Burr and S. Loo, On Rado numbers II, unpublished.

% \end{thebibliography}

\end{document}  



